\subsection{Automorphic Forms and Fermion Masses --- arXiv:2010.07952}

\textbf{Authors:} Gui-Jun Ding, Ferruccio Feruglio, Xiang-Gan Liu

\paragraph{主な主張}
複数モジュライをもつ超対称理論にモジュラー不変性を一般化し、$Sp(2g,\mathbb{Z})$ に対するSiegelモジュラー形式をYukawa結合とするフェルミオン質量模型を構成できることを示した\cite{Ding:2020zxw}。

\paragraph{新規性}
従来の一変数のモジュラー形式を、対称空間 $G/K$ と離散群 $\Gamma$ に基づく保型形式へ拡張し、種数2の具体的なクォーク・レプトン模型を与えた。

\paragraph{位置付け}
多モジュライのモジュラーフレーバー模型を組み立てるための一般的なボトムアップの枠組みである。

\paragraph{コメント（読書メモ）}
読書メモ：モジュライ空間を、非コンパクト連続群 $G$ とその最大コンパクト部分群 $K$ による対称空間 $G/K$ を離散群 $\Gamma$ で割ったものとして捉える。物質場が $\Gamma$ の非自明な表現で変換されるため、フレーバー対称性の起源をこの離散双対群に求めるのが魅力的に思える。Feruglioによる $SL(2,\mathbb{Z})$ の模型構築\cite{Feruglio:2017spp}は $G=SL(2,\mathbb{R})$, $K=SO(2)$ という選択に対応し、本論文はその一般化を行うボトムアップの研究として読む。
